\documentclass[11pt,a4paper]{article}
\usepackage[T1]{fontenc}
\usepackage[utf8]{inputenc}
\usepackage{amsmath,amssymb,graphicx,hyperref,longtable,booktabs,array}
\usepackage[margin=2.2cm]{geometry}
\providecommand{\tightlist}{\setlength{\itemsep}{0pt}\setlength{\parskip}{0pt}}
\setlength{\parindent}{0pt}\setlength{\parskip}{6pt}

\begin{document}
\hypertarget{piece-1-verification-of-the-chat-1-prototype-30-oct-2026}{%
\section{Piece 1: verification of the chat-1 prototype (30 Oct
2026)}\label{piece-1-verification-of-the-chat-1-prototype-30-oct-2026}}

Scope: check the inherited simglucose + EKF prototype
(\texttt{prototype/}) before building on it. The prototype was written
by Claude (AI assistant) in an earlier chat; nobody on the team has yet
run or reviewed it, so it must not be presented as team work until Piece
3-4 re-implement or sign off each part.

\hypertarget{checked-and-confirmed}{%
\subsection{Checked and confirmed}\label{checked-and-confirmed}}

\begin{itemize}
\tightlist
\item
  Raw-CGM RMSE vs true plasma recomputed directly from \texttt{sim.npz}
  on the 13 held-out patients: 8.71 (NS=0), 9.46 (NS=0.3), 15.08 (NS=1).
  Matches the handoff.
\item
  Every row of the handoff results table, hypoglycaemia delays and
  tau-recovery figures matches the saved results JSON files
  (tau-recovery correlation is 0.0 to 0.36, not "about 0.1 to 0.2").
\item
  Patient exclusions (child\#002, \#004, \#008), 14 tune / 13 test split
  and tau range (6.6 to 19.9 min, mean 11.9, sd 3.2) reproduce.
\item
  simglucose 0.2.11 source: subcutaneous state obeys dx12/dt =
  -ksc\emph{x12 + ksc}x3, so the simulator lag is exactly first-order
  with tau = 1/ksc.
\item
  CGM samples update at minute indices = 2 (mod 3), which matches the
  \texttt{arange(2,\ N,\ 3)} sampling in \texttt{evaluate.py}; no
  staleness offset.
\end{itemize}

\hypertarget{resolved-open-item}{%
\subsection{Resolved open item}\label{resolved-open-item}}

\begin{itemize}
\tightlist
\item
  Handoff said the 0.98 lag-1 autocorrelation of sensor noise was "not
  verified". Verified: simglucose \texttt{noise\_gen.py} draws
  Johnson-transformed noise every 15 min (PACF 0.7) and
  cubic-interpolates it to the sample time. Measured noise
  autocorrelation: 0.98 at 3 min, 0.72 at 15 min, 0.47 at 30 min, 0.24
  at 60 min. So the simulator noise is slow and drift-like with no fast
  component; "noise dominates at NS=1" is partly a property of this
  noise model, not of real CGM.
\end{itemize}

\hypertarget{issues-to-fix-carried-into-later-pieces}{%
\subsection{Issues to fix (carried into later
pieces)}\label{issues-to-fix-carried-into-later-pieces}}

\begin{enumerate}
\def\labelenumi{\arabic{enumi}.}
\tightlist
\item
  Noise model has no fast component and is unrealistic for real sensors
  (Piece 2).
\item
  Several tuned process-noise values sit at the grid edge (qr = 0.3 at
  NS=0 for every EKF variant); widen the grid (Piece 3).
\item
  Noise sd and phi are estimated from (cgm - true subcutaneous glucose),
  which needs ground truth; fine in simulation, not available on real
  data.
\item
  At NS=0 online-tau EKF (RMSE 5.05) is worse than fixed tau=10 (2.36),
  and tau error (57-63\%) is worse than using the 10-min prior (22\%).
  Treat as an open problem and a tuning question, not a settled result.
\item
  Lag is exactly first-order in the generator, so all results are
  best-case (no mismatch yet).
\item
  \texttt{fc30\_*} forecast metrics are not meaningful; do not report.
\item
  Time below 70 mg/dL on the test set is 4.6\% (handoff quotes about
  6.6\% overall); minor.
\end{enumerate}

\hypertarget{differences-from-the-project-plan-doc}{%
\subsection{Differences from the project plan
doc}\label{differences-from-the-project-plan-doc}}

\begin{itemize}
\tightlist
\item
  Prototype samples at 3 min (Dexcom model in simglucose), plan assumes
  5 min; run 1, 5 and 15 min sensitivity in Piece 3.
\item
  Prototype uses a nonlinear EKF with ln(tau) in the state;
  plan\textquotesingle s core method is a linear Kalman filter and RTS
  smoother with a regularised-inversion cross-check. Keep the EKF as a
  comparison.
\item
  Plan\textquotesingle s hypotheses H1 to H4 are not yet tested on
  mismatched generators (G2, G3).
\end{itemize}

\hypertarget{proposed-pieces}{%
\subsection{Proposed pieces}\label{proposed-pieces}}

\begin{enumerate}
\def\labelenumi{\arabic{enumi}.}
\tightlist
\item
  Verify prototype (this note).
\item
  Realistic sensor layer and mismatch generators (fast noise, extra
  delay, drift, compression lows).
\item
  Evaluation harness: baselines, widened tuning, subject-wise splits,
  bootstrap intervals, analytic ramp test.
\item
  Linear Kalman, RTS smoother and regularised inversion; cross-check.
\item
  Identifiability study of tau (profile likelihood, population prior,
  fingerstick-timed regression).
\item
  Meal and insulin-aware prior (optional).
\item
  OhioT1DM access and exploratory real-data check.
\item
  Ablations, final results, report.
\end{enumerate}
\end{document}
